% =====================================================================
%  CARNET D'ENTRAINEMENT — FICHE 1 : Les chiffres significatifs
%  THÈME 4 — Logarithmes, pH & incertitude (justesse / fidélité)
%  Fiche TUTO
% =====================================================================
\documentclass[11pt,a4paper]{article}
\usepackage[utf8]{inputenc}
\usepackage[T1]{fontenc}
\usepackage{lmodern}
\usepackage[french]{babel}
\usepackage{amsmath,amssymb,mathtools}
\usepackage{icomma}
\usepackage{siunitx}
\sisetup{output-decimal-marker={,},per-mode=symbol}
\usepackage[version=4]{mhchem}
\usepackage{xcolor}
\usepackage{tikz}
\usepackage[most]{tcolorbox}
\usepackage{enumitem}
\usepackage{booktabs}
\usepackage{fancyhdr}
\usepackage[a4paper,margin=2cm,headheight=26pt,headsep=8pt]{geometry}

\definecolor{primary}{RGB}{150,98,162}
\definecolor{primarydark}{RGB}{92,52,110}
\definecolor{formulacol}{RGB}{198,93,9}
\definecolor{reglecol}{RGB}{150,98,162}
\definecolor{defcol}{RGB}{0,114,178}
\definecolor{excol}{RGB}{0,135,90}
\definecolor{attncol}{RGB}{200,40,55}
\newcommand{\cs}{chiffre significatif}
\newcommand{\CS}{chiffres significatifs}

\pagestyle{fancy}\fancyhf{}
\fancyhead[L]{\small\textcolor{primarydark}{\textbf{Fiche 1} — Chiffres significatifs}}
\fancyhead[R]{\small\textcolor{primarydark}{\textsc{Tuto · Thème 4}}}
\fancyfoot[C]{\small\thepage}
\renewcommand{\headrulewidth}{0.8pt}
\renewcommand{\headrule}{\hbox to\headwidth{\color{primary}\leaders\hrule height \headrulewidth\hfill}}

\tcbset{boxbase/.style={enhanced,breakable,boxrule=0.9pt,arc=2.5pt,
  left=3mm,right=3mm,top=1.6mm,bottom=1.6mm,fonttitle=\bfseries,coltitle=white}}
\newtcolorbox{regle}[1][]{boxbase,colback=reglecol!7,colframe=reglecol,
  title={\textsc{Règle}\ifx\relax#1\relax\relax\else\ ---\ \textmd{\itshape #1}\fi}}
\newtcolorbox{formule}[1][]{boxbase,colback=formulacol!7,colframe=formulacol,
  title={\textsc{Formule}\ifx\relax#1\relax\relax\else\ ---\ \textmd{\itshape #1}\fi}}
\newtcolorbox{defi}[1][]{boxbase,colback=defcol!6,colframe=defcol,
  title={\textsc{Définition}\ifx\relax#1\relax\relax\else\ ---\ \textmd{\itshape #1}\fi}}
\newtcolorbox{exemple}[1][]{boxbase,colback=excol!6,colframe=excol,
  title={\textsc{Exemple}\ifx\relax#1\relax\relax\else\ : \textmd{\itshape #1}\fi}}
\newtcolorbox{attention}[1][]{boxbase,colback=attncol!6,colframe=attncol,
  title={\textsc{Attention}\ifx\relax#1\relax\relax\else\ : \textmd{\itshape #1}\fi}}

\setlength{\parindent}{0pt}\setlength{\parskip}{2pt}

\begin{document}

\begin{center}
{\LARGE\bfseries\color{primarydark} Thème 4 — Logarithmes, pH \& incertitude}\\[1pt]
{\large\color{primary} La précision, du calcul de pH à l'épaisseur du doute}
\end{center}
\vspace{1mm}

Ce thème regroupe deux usages « précision » propres à la chimie : la règle des \CS\ dans un
\textbf{logarithme} (indispensable pour le pH), et le vocabulaire de la qualité d'une mesure :
\textbf{justesse}, \textbf{fidélité} et \textbf{incertitude}.

\begin{formule}[le pH et sa règle de \CS]
\[ \boxed{\;\mathrm{pH}=-\log_{10}\bigl([\ce{H3O+}]\bigr)\;} \qquad
   \boxed{\;[\ce{H3O+}]=10^{-\mathrm{pH}}\;} \]
Dans un logarithme, \emph{seule la partie décimale} (après la virgule) est significative ; la partie
entière ne fait que rappeler la puissance de dix. D'où la règle :
\[ \text{nombre de \emph{décimales} du pH} \;=\; \text{nombre de \CS\ de } [\ce{H3O+}]. \]
\end{formule}

\begin{exemple}[dans les deux sens]
$[\ce{H3O+}]=\num{2,5e-4}\,\si{\mol\per\litre}$ ($2$~CS) $\Rightarrow$ pH à \textbf{2 décimales} :
$\mathrm{pH}=-\log(\num{2,5e-4})=\num{3,60}$. Les décimales « $60$ » sont les seuls chiffres
significatifs ; le « $3$ » situe l'ordre de grandeur ($10^{-4}$). \\
Inversement, $\mathrm{pH}=\num{2,00}$ ($2$ décimales) $\Rightarrow [\ce{H3O+}]=10^{-2,00}=\num{1,0e-2}\,\si{\mol\per\litre}$
($2$~CS), et non $\num{1e-2}$.
\end{exemple}

\begin{defi}[justesse et fidélité]
\textbf{Justesse} : proximité de la mesure à la \emph{vraie} valeur (mesure non biaisée).
\textbf{Fidélité} : proximité des mesures \emph{entre elles} quand on répète (reproductibilité). Le
nombre de \CS\ traduit la \emph{fidélité}, jamais la justesse : une balance déréglée peut afficher
$4$~CS tout en étant fausse.
\end{defi}

\begin{center}
\begin{tikzpicture}[scale=0.52,
   cible/.style={fill=primary},
   anneaux/.pic={\fill[red!12] (0,0) circle (1.5);\fill[white] (0,0) circle (1.0);
     \fill[red!12] (0,0) circle (0.5);\draw[black!55] (0,0) circle (1.5);
     \draw[black!55] (0,0) circle (1.0);\draw[black!55] (0,0) circle (0.5);
     \fill[black!70] (0,0) circle (0.06);}]
  \begin{scope}[xshift=0cm]\pic{anneaux};
     \fill[cible](0,0)circle(0.1);\fill[cible](0.2,0.13)circle(0.1);
     \fill[cible](-0.13,0.17)circle(0.1);\fill[cible](0.11,-0.16)circle(0.1);
     \node[below,align=center,font=\scriptsize] at (0,-1.95){juste\\ \& fidèle};\end{scope}
  \begin{scope}[xshift=4.4cm]\pic{anneaux};
     \fill[cible](1.0,0.9)circle(0.1);\fill[cible](1.2,1.05)circle(0.1);
     \fill[cible](0.86,1.0)circle(0.1);\fill[cible](1.05,0.78)circle(0.1);
     \node[below,align=center,font=\scriptsize] at (0,-1.95){fidèle\\ non juste};\end{scope}
  \begin{scope}[xshift=8.8cm]\pic{anneaux};
     \fill[cible](-0.9,0.3)circle(0.1);\fill[cible](0.8,-0.6)circle(0.1);
     \fill[cible](0.2,1.0)circle(0.1);\fill[cible](-0.4,-0.9)circle(0.1);
     \node[below,align=center,font=\scriptsize] at (0,-1.95){juste\\ non fidèle};\end{scope}
  \begin{scope}[xshift=13.2cm]\pic{anneaux};
     \fill[cible](0.9,1.0)circle(0.1);\fill[cible](1.3,0.6)circle(0.1);
     \fill[cible](-1.1,-0.4)circle(0.1);\fill[cible](0.4,-1.2)circle(0.1);
     \node[below,align=center,font=\scriptsize] at (0,-1.95){ni juste\\ ni fidèle};\end{scope}
\end{tikzpicture}
\end{center}
\vspace{-1mm}
{\footnotesize Le centre = vraie valeur ; chaque point = une mesure répétée. \textbf{Fidélité} = points serrés ; \textbf{justesse} = points centrés.}

\begin{regle}[incertitude absolue et relative]
Le résultat s'écrit $x = x_{\text{mes}} \pm \Delta x$. L'\textbf{incertitude absolue} $\Delta x$ (même
unité que $x$) est fixée par le dernier \cs\ écrit : $m=\qty{13,02}{\gram}$ sous-entend
$\Delta m=\qty{0,01}{\gram}$. L'\textbf{incertitude relative} $\dfrac{\Delta x}{x}$ (sans unité,
souvent en~\%) permet de comparer des mesures. À $n$~\CS, elle est de l'ordre de $10^{-(n-1)}$ :
$2$~CS~$\sim 1\%$, $3$~CS~$\sim0{,}1\%$, $4$~CS~$\sim0{,}01\%$.
\end{regle}

\end{document}
